% ===========================================================
%  Praeambel -- Pakete und Einstellungen
% ===========================================================

% ---------- Sprache & Kodierung ----------
\usepackage[T1]{fontenc}
\usepackage[utf8]{inputenc}
\usepackage[ngerman]{babel}      % Englisch: [english]{babel}
\usepackage{csquotes}           % von biblatex benoetigt

% ---------- Schrift & Mikrotypografie ----------
\usepackage{lmodern}
\usepackage{microtype}

% ---------- Mathematik ----------
\usepackage{amsmath,amssymb,mathtools}
\usepackage{bm}

% ---------- Einheiten & Chemie ----------
\usepackage{siunitx}
\sisetup{
    locale = DE,                % Englisch: locale = US
    per-mode = symbol,
    range-units = single,
}
\usepackage[version=4]{mhchem}

% ---------- Abbildungen & Tabellen ----------
\usepackage{graphicx}
\graphicspath{{abbildungen/}}
\usepackage{booktabs}
\usepackage{longtable}
\usepackage{multirow}
\usepackage{float}
\usepackage[labelfont=bf,font=small]{caption}
\usepackage{subcaption}

% ---------- Code-Listings (fuer Methodenteil) ----------
\usepackage{listings}
\lstset{
    basicstyle=\ttfamily\footnotesize,
    breaklines=true,
    frame=single,
    numbers=left,
    numberstyle=\tiny,
    captionpos=b,
    showstringspaces=false,
    language=Python,
}

% ---------- Literatur (Zotero + Better BibTeX) ----------
\usepackage[
    backend=biber,
    style=numeric-comp,         % Alternative: authoryear
    sorting=none,               % Nummerierung nach Erstzitat
    maxbibnames=99,
    giveninits=true,
    doi=true,
    url=false,
    isbn=false,
]{biblatex}
\addbibresource{literatur/literatur.bib}

% ---------- Querverweise (immer zuletzt laden) ----------
\usepackage[
    hidelinks,
    breaklinks=true,
    pdfusetitle,
]{hyperref}
\usepackage[noabbrev,nameinlink]{cleveref}

% ---------- Eigene Makros ----------
\newcommand{\amp}{AMP}
\newcommand{\esm}{ESM-2}

% ---------- Platzhalter sichtbar machen ----------
\usepackage{xcolor}
\newcommand{\TODO}[1]{\textcolor{red}{\textbf{[TODO:} #1\textbf{]}}}
